\documentclass[10pt]{article}
\usepackage[margin=1in]{geometry}
\usepackage[T1]{fontenc}
\usepackage{lmodern,microtype,amsmath,amssymb,amsthm}
\theoremstyle{definition}
\newtheorem{assumption}{Assumption}
\begin{document}
\section*{Convex Optimization with Local Linear Access}
We consider minimizing $f:E\to\mathbb{R}$ over a closed convex set $C$ in a real Hilbert space $E$. On $C$, the objective is differentiable and convex, with quadratic upper models of hidden curvature $L\geq0$. A given $x^0\in C$ is within distance $R\geq0$ of a minimizer $x^\star$, with hidden initial gap bound $\Delta\geq0$. Each oracle call returns either an exact value and gradient or a linear minimizer over a ball intersected with $C$. Given $\varepsilon>0$, the required output $\bar{x}$ belongs to $C$ and satisfies $f(\bar{x})-f(x^\star)\leq\varepsilon$ almost surely. Only $(E,R,\varepsilon,x^0)$ are public inputs.

The optimization problem is
\[
\min_{x\in C} f(x),
\]
where the norm $\|\cdot\|$ is induced by the inner product $\langle\cdot,\cdot\rangle$ on $E$. The admissible instances and oracle replies satisfy the following conditions.

\begin{assumption}[Feasible geometry]
The set $C\subseteq E$ is closed and convex, and $x^0\in C$.
\end{assumption}

\begin{assumption}[First-order models]
The function $f$ is Fr\'echet differentiable at every point of $C$, with gradient $\nabla f(x)\in E$. For every $x,y\in C$,
\[
\begin{aligned}
f(x)+\langle\nabla f(x),y-x\rangle &\leq f(y),\\
f(y) &\leq f(x)+\langle\nabla f(x),y-x\rangle
       +\frac{L}{2}\|y-x\|^2.
\end{aligned}
\]
\end{assumption}

\begin{assumption}[Attainment and initialization]
There exists $x^\star\in C$ such that
\[
\begin{aligned}
f(x^\star) &\leq f(x) \qquad \text{for every }x\in C,\\
\|x^0-x^\star\| &\leq R,\\
f(x^0)-f(x^\star) &\leq\Delta.
\end{aligned}
\]
\end{assumption}

\begin{assumption}[Exact value and local linear access]
A single oracle, with ID $0$, accepts two types of queries. A value--gradient query at $x\in C$ returns the pair $(f(x),\nabla f(x))$. A local linear query $(x,r,g)$, with $x\in C$, $r\geq0$, and $g\in E$, returns a point $v\in E$ satisfying
\[
\begin{aligned}
v&\in C,\qquad \|v-x\|\leq r,\\
\langle g,v\rangle&\leq\langle g,z\rangle
\quad\text{for every }z\in C\text{ with }\|z-x\|\leq r.
\end{aligned}
\]
These guarantees hold for every query history and every realization of oracle randomness. The local linear minimizer may depend on both. Each query of either type costs one call to the same oracle. Internal algorithm randomness is independent of oracle randomness.
\end{assumption}

The algorithm is selected using only the public inputs $(E,R,\varepsilon,x^0)$; neither $L$, $\Delta$, nor the optimal value $f(x^\star)$ is supplied. The objective and feasible set are accessed through the oracle. Queries may be adaptive. The query interface accepts points throughout $E$ and arbitrary real radii, but replies outside the domains specified above are unrestricted. An externally specified budget $K\in\mathbb{N}$ may depend on the full parameter tuple $(E,R,\varepsilon,x^0,L,\Delta)$ and counts both query types together.

The output $\bar{x}\in E$ is returned after at most $K$ oracle calls. Writing $\mathbb{P}$ for the joint law of algorithm and oracle randomness, success requires a minimizer $x^\star\in C$ such that
\[
\bar{x}\in C
\quad\text{and}\quad
f(\bar{x})-f(x^\star)\leq\varepsilon
\qquad \mathbb{P}\text{-almost surely}.
\]
The minimizer witnessing the initialization bounds may be used here, since all minimizers on $C$ have the same objective value.
\end{document}
